\subsection{Equilibrio en el análisis de ingreso nacional}

El estudio de análisis estático tiene aplicaciones en áreas de la economía, como por ejemplo, el del modelo de ingreso nacional keynesiano más simple, donde
\begin{equation}
\label{ecuacion_equilibrio_ingreso_nacional_simple_1}
\begin{aligned}
Y &= C + I_0 + G_0 \\
C &= a + bY
\end{aligned}
\end{equation}
En el cual \(a > 0, 0<b<1\). En esta ecuación \(Y\) y \(C\) representan las variables endógenas de ingreso nacional y gasto de consumo (planificado), respectivamente, \(I_0\) representa la inversion y \(G_0\) los gastos de gobierno determinados de manera exógena. 

La primera ecuación es una condición de equilibrio \textit{(ingreso nacional = gasto total planeado)}, mientras que la segunda función es de comportamiento. La segunda función, de consumo, cosnta de dos parametros de consumo \(a\) y \(b\), los cuales representan el gasto de consumo autónomo y la propensión marginal al consumo.

Estas dos ecuaciones, con variables endógenas, no son dependientes ni inconsistentes. Se podrían hallar valores de equilibrio de ingreso y gasto de consumo \(Y^\ast\) y \(C^\ast\) en términos de los parámetros \(a\) y \(b\) y las variables exógenas \(I_0\) y \(G_0\).

Al sustituir la segunda ecuación en la primera reduce la ecuación \ref{ecuacion_equilibrio_ingreso_nacional_simple_1} a una sola ecuación en la variable \(Y\)
$$
\begin{aligned}
Y &= a + bY + I_0 + G_0  \qquad \text{O bien}\\
(1-b)Y &= a + I_0 + G_0 \qquad (\text{Reuniendo los terminos relacionados con } Y)
\end{aligned}
$$
Para hallar el valor de solución de \(Y\) (ingreso nacional de equilibrio), sólo se debe dividir entre \((1-b)\)
\[
Y^\ast=\frac{a + I_0 + G_0}{1-b}
\]
El valor de solución se expresa por completo en término de parámetros y variables exógenas, los cuales son los datos del modelo. Si colocamos el valor de solución \(Y^*\) en \(C = a + bY\), se produce entonces el nivel de equilibrio de gasto de consumo, donde 
$$
\begin{aligned}
C^\ast &= a + \frac{b(a + I_0 + G_0)}{1-b}\\
C^\ast &= \frac{a(1-b)+b(a+I_0+G_0)}{1-b}=\frac{a+b(I_0+G_0)}{1-b}
\end{aligned}
$$
Como \(Y^\ast\) y \(C^\ast\) tienen la expresión \((1-b)\) en el denominador, es necesario restringir que \(b \neq 1\) para evitar la división entre cero. Esto se realiza bajo el supuesto que \(b\), la propensión marginal al consumo, es una fracción positiva. Para que \(Y^\ast\) y \(C^\ast\) sean positivos, los numeradores de sus ecuaciones tambien deben serlo. Los gasots exógenos como \(I_0\) y \(G_0\) son positivos, al igual que el parámetro \(a\).

Este modelo es de extrema simplicidad y crudeza, pero se pueden construir otros modelos de determinación de ingreso nacional con distintos grados de complejidad.

\begin{ejercicio}
    \begin{enumerate}
        \item Dado el siguiente modelo 
        $$
        \begin{aligned}
            Y &= C+I_0+G_0 \\
            C &=a+b(Y-T) \qquad &&T\text{ es impuestos} \\
            T &= d+tY \qquad &&t\text{ tasa de impuestos sobre la renta}
        \end{aligned}
        $$
        \((a)\) ¿Cuántas variables endógenas hay?
        \((b)\) Determine \(Y^\ast,T^\ast\) y \(C^\ast\).
        \item Sea el modelo de ingreso nacional 
        $$
        \begin{aligned}
            Y &= C+I_0+G_0 \\
            C &=a+b(Y-T_0)  \\
            G &= gY
        \end{aligned}
        $$
        \((a)\) Identifique las variables endógenas.
        \((b)\) Dé el significado económico del parámetro \(g\).
        \((c)\) Determine el ingreso nacional de equilibrio.
        \((d)\) ¿Qué restricción se requiere en los parámetros para que exista una solución?
        \item Determine \(Y^\ast\) y \(C^\ast\) a partir de lo siguiente
        $$
        \begin{aligned}
            Y &= C+I_0+G_0 \\
            C &=35+6Y^{1/2}  \\
            I_0 &= 16 \\
            G_0 & = 14
        \end{aligned}
        $$
    \end{enumerate}
    \emph{Respuestas}
    \begin{enumerate}
        \item Hay 3 variables endógenas \(Y,C\) y \(T\)
        $$
        \begin{aligned}
            Y^\ast &= \frac{a-bd+I_0+G_0}{1-b+bt} \\
            C^\ast &= \frac{a-bd+b(1-t)(I_0G_0)}{1-b+bt}  \\
            T^\ast &= \frac{d-db+t(a+I_0+G_0)}{1-b+bt}
        \end{aligned}
        $$
        Para entender que es cada variable ver cuadro \ref{Significado_Variables_ingreso_nacional}
        \item \((a)\) Las variables endógenas en este modelo son el ingreso nacional \((Y)\), el consumo \((c)\) y el gasto público \((G)\).\\
        \((b)\) \(g\) representa la propension marginal al gasto público. Indica cuanto aumenta el gasto del gobierno por cada unidad de ingreso nacional.\\
        \((c)\) \[
        Y^\ast=\frac{a-bT_0+I_0}{1-b-g}
        \]
        \((d)\) \(1-b-g > 0\) y \(b+g < 1\) Ya que no se puede dividir por \(0\) y si el denominador es negativo nos daria un resultado negativo.
        \item El procedimiento para el resultado de\(Y^ast\)
        $$
        \begin{aligned}
            & Y = C+16+14 \\
            & Y = 35+6Y^{1/2}+30 \\
            & Y = 6Y^{1/2}+65 \\
            & Y-65 = 6Y^{1/2} \\
            & (Y-65)^2=(6Y^{1/2})^2 \\ 
            & Y^2 - 2\times 65 \times Y + 65^2 = 36Y \\
            & Y^2 -130Y+65^2=36 Y \\
            & Y^2 -166Y + 65^2 = 0 \\
            & Y_1 = 83-6\sqrt{74} \text{(No verifica, es negativo)} \\
            & Y_2^\ast = 83+6\sqrt{74}\approx 134,6
        \end{aligned}
        $$
        Para obtener \(C\) remplazo \(Y\)
        $$
        \begin{aligned}
            & C = 35+6\sqrt{Y_2^\ast} 
        \end{aligned}
        $$
    \end{enumerate}
\end{ejercicio}
        
\begin{table}[h]
\centering
\begin{tabular}{|c|l|l|}
\hline
\textbf{Símbolo} & \textbf{Nombre} & \textbf{Significado Económico} \\ \hline
$Y$ & Ingreso Nacional (PIB) & Valor total de la producción final. \\ \hline
$C$ & Consumo & Gasto total de los hogares. \\ \hline
$T$ & Impuestos Totales & Recaudación fiscal del Estado. \\ \hline
$a$ & Consumo Autónomo & Consumo que no depende del ingreso. \\ \hline
$b$ & PMC & Propensión Marginal a Consumir ($0 < b < 1$). \\ \hline
$d$ & Impuestos Autónomos & Impuestos fijos, independientes del ingreso. \\ \hline
$t$ & Tasa Impositiva & Impuesto proporcional sobre la renta. \\ \hline
$I_0$ & Inversión Exógena & Inversión planeada por las empresas. \\ \hline
$G_0$ & Gasto Público & Gasto realizado por el gobierno. \\ \hline
$g$ & Propersión al gasto & Incremento del gasto por cada unidad de \(Y\) \\ \hline
\end{tabular}
\caption{Tabla de variables}
\label{Significado_Variables_ingreso_nacional}
\end{table}
